% Figure: pressure-enthalpy (p-h) diagram for a refrigerant, the canonical
% representation of a vapour-compression cycle. The dome is the saturation
% boundary; the cycle is the four-leg loop from the worked R-134a example:
% 1 saturated vapour at low pressure, 2 after isentropic compression,
% 3 saturated liquid at high pressure, 4 after isenthalpic throttling.
% Pressure axis is logarithmic in practice; here drawn schematically.
\begin{figure}[htbp]
\centering
\definecolor{engteal}{RGB}{30,120,120}
\begin{tikzpicture}
\begin{axis}[
  width=12cm, height=8cm,
  xlabel={specific enthalpy $h$ (kJ/kg)},
  ylabel={pressure $p$ (kPa, log scale)},
  ymode=log,
  xmin=80, xmax=300, ymin=100, ymax=1300,
  axis lines=left,
  log ticks with fixed point,
  tick label style={font=\scriptsize},
  label style={font=\small},
  clip=true,
]
% saturation dome: left limb (saturated liquid, h_f rises with p),
% right limb (saturated vapour, h_g) meeting near the top (critical region)
\addplot[engblue, very thick, smooth] coordinates {
  (95,120) (100,250) (107,500) (118,800) (135,1100) (165,1250)
};
\addplot[only marks, mark=*, engblack, mark size=1.6pt] coordinates {(165,1250)};
\addplot[engblue, very thick, smooth] coordinates {
  (165,1250) (210,1100) (235,800) (244,500) (248,250) (250,120)
};
\node[engblue, font=\scriptsize] at (axis cs:175,300) {two-phase};
% vapour-compression cycle (R-134a worked example, schematic):
% 1: saturated vapour, low p ~ 133 kPa, h1 = 235.3
% 2: after isentropic compression, p = 1000 kPa, h2 = 277 (superheated, right of dome)
% 3: saturated liquid, p = 1000 kPa, h3 = 107.3
% 4: after throttle to 133 kPa, h4 = 107.3, two-phase
\addplot[only marks, mark=*, engred, mark size=2.0pt] coordinates {
  (235.3,133) (277,1000) (107.3,1000) (107.3,133)
};
% 1 -> 2 compression
\addplot[engred, very thick] coordinates {(235.3,133) (277,1000)};
% 2 -> 3 condensation (constant p, leftwards along high isobar)
\addplot[engorange, very thick] coordinates {(277,1000) (107.3,1000)};
% 3 -> 4 throttle (constant h, downwards)
\addplot[engred, very thick] coordinates {(107.3,1000) (107.3,133)};
% 4 -> 1 evaporation (constant p, rightwards along low isobar)
\addplot[engteal, very thick] coordinates {(107.3,133) (235.3,133)};
\node[engblack, font=\scriptsize, anchor=north] at (axis cs:235.3,128) {1};
\node[engblack, font=\scriptsize, anchor=west]  at (axis cs:279,1000) {2};
\node[engblack, font=\scriptsize, anchor=east]  at (axis cs:105,1000) {3};
\node[engblack, font=\scriptsize, anchor=north] at (axis cs:107.3,128) {4};
\node[engred, font=\scriptsize, anchor=south west] at (axis cs:250,560) {compressor};
\node[engorange, font=\scriptsize, anchor=south] at (axis cs:192,1010) {condenser};
\node[engteal, font=\scriptsize, anchor=north] at (axis cs:171,128) {evaporator};
\end{axis}
\end{tikzpicture}
\caption[Pressure-enthalpy diagram of a vapour-compression cycle]{The
pressure-enthalpy diagram for a refrigerant, with the saturation dome
(blue) and the four-leg vapour-compression cycle of the worked R-134a
example inscribed. The compressor raises the saturated vapour (state~1)
to the condenser pressure (state~2), a near-vertical line on the
log-pressure axis; the condenser rejects heat at constant pressure to the
saturated liquid (state~3); the throttle drops the pressure at constant
enthalpy into the two-phase region (state~4); and the evaporator absorbs
heat at constant pressure back to state~1. The horizontal width of the
evaporator leg is the specific cooling capacity $q_L=h_1-h_4$; the width
of the compressor leg is the specific work $w_c=h_2-h_1$. Their ratio is
the coefficient of performance. The pressure axis is drawn schematically;
working charts use a logarithmic pressure scale.}
\label{fig:vol04ch02:ph-refrigerant-cycle}
\end{figure}
